\begin{tikzpicture}[
    tensor/.style={rectangle, draw, fill=white, text width=2em, text centered, minimum height=2em},
    % % --- NUEVOS ESTILOS PARA ETIQUETAS ---
    lbl_fwd/.style={above, pos=0.5, font=\footnotesize}, % 'pos=0.5' centra la etiqueta
    lbl_bwd/.style={below, pos=0.5, font=\footnotesize},
    lbl_left/.style={left, pos=0.5, font=\footnotesize},
    lbl_right/.style={right, pos=0.5, font=\footnotesize},
    layer_function/.style={rectangle, draw, fill=white, text width=4em, text centered, rounded corners, minimum height=2em},
    operation/.style={circle, draw, fill=white, minimum size=0.1em, font=\tiny},
    arrow/.style={-latex, thick}
]

\tikzset{node distance=1.75cm and 1.5cm}

%%%% FORWARD %%%%
% draw linear block z()
\node[layer_function] (z) at (0,0) {$z()$};

% --- MODIFIED: Changed from (-2,0) to (-3,0) ---
\draw[arrow] ($(z)+(-3,0)$) -- node[lbl_fwd] {$a_{m-1}$} (z.west);

% draw activation layer_function phi()
\node[layer_function, right=of z] (phi) {$\phi()$};

% draw output activation a_m (Already 3cm, no change needed)
\draw[arrow] (phi.east) -- node[lbl_fwd, pos=0.75] {$a_{m}$} ($(phi)+(4,0)$);

% draw arrow from z to phi
% annotate it as $z_m$
\draw[arrow] (z.east) -- node[lbl_fwd] {$z_{m}$} (phi.west);

% draw theta parameters
\node[tensor, below=of z] (theta) {$\theta$};

% theta to z
\draw[arrow] (theta) -- (z);


% % --- DEFINE PARAMETER UPDATE CHAIN FIRST --- % %
\def\updateopspacing{0.5cm}
\node[operation, below=\updateopspacing of theta] (mult_eta) {$\times$};
\node[operation, below=\updateopspacing of mult_eta] (negative) {$-$};
\node[left=0.5cm of mult_eta] (eta) {$\eta$};


%%%% BACKWARD (NOW POSITIONED RELATIVE TO UPDATE CHAIN) %%%%

% % draw z gradient, positioned 0.75cm below 'negative'
\node[layer_function, below=0.75cm of negative] (dz) {$\nabla_z$};

% % draw phi gradiente
\node[layer_function, right=of dz] (dphi) {$\nabla_\phi$};

% Theres a multication between upstream and dphi
\node[operation] (mult) at ($(dphi)+(2,0)$) {$\times$};

% --- MODIFIED: Changed from (1,0) to (3,0) and simplified label ---
\draw[arrow] ($(mult)+(2,0)$) -- (mult.east) node[lbl_bwd] {$\frac{\partial \mathcal{L}}{\partial a_{m}}$};
\draw[arrow] (mult.west) -- (dphi.east);

% draw a circle with $\times in the midle of dz and dphi
\node[operation] (mult2) at ($(dz)!.5!(dphi)$) {$\times$};
\draw[arrow] (dphi.west) -- (mult2.east);
\draw[arrow] (mult2.west) -- (dz.east);

% --- MODIFIED: Changed from (-2,0) to (-3,0) ---
\draw[arrow] (dz.west) -- node[lbl_bwd] {$\frac{\partial \mathcal{L}}{\partial a_{m-1}}$} ($(dz)+(-3,0)$);


% % --- CONNECT ALL THE NODES --- % %
\draw[arrow] (eta.east) -- (mult_eta.west);
\draw[arrow] (mult_eta.north) -- (theta.south);
\draw[arrow] (negative.north) -- (mult_eta.south);

% --- MODIFIED: Added name (dtheta) to the label node ---
\draw[arrow] (dz.north) -- node[lbl_left] (dtheta) {$\frac{\partial \mathcal{L}}{\partial \theta}$} (negative.south);

% % --- BACKGROUND LAYER BOX --- % %
\begin{pgfonlayer}{background}
    \node [
        draw,
        dotted,
        thick,
        orange,
        rounded corners,
        % --- MODIFIED: Added (dtheta) to fit list ---
        fit=(z) (phi) (theta) (mult_eta) (negative) (dz) (dphi) (mult) (mult2) (dtheta),
        % --- MODIFIED: Reduced inner sep for better spacing ---
        inner sep=0.35cm
    ] (layer_box) {};
    % Add horizontal line through the box at theta level (background so it sits behind nodes)
    \draw [gray, dashed] ($(theta)+(-2.5,0)$) -- ($(theta)+(7,0)$);
\end{pgfonlayer}

% --- Placed your label relative to the box's corner ---
\node[font=\footnotesize, orange, above right, xshift=2mm] at (layer_box.north east) {$L_m$};


\end{tikzpicture}
